\subsection{FPTRANSA Accuracy Discovery}
The (:ref:FPTRANSA.cpuid.EXTENSIONS.DIRECTORY.FEATURES.FPTRANSA:) predicate requires (:ref:FP.cpuid.EXTENSIONS.DIRECTORY.FEATURES.FP:) and reports that at least one approximate transcendental
contract may be available. Software determines each operation\textquotesingle{}s availability separately from its
accuracy-contract result.
Class \texttt{0x00000001}, leaf \texttt{0x0002} uses indexes 1 through \texttt{0x0044} as a sparse space of stable
16-bit accuracy-contract IDs. Index zero is the common (:term:base.terms.header:) and reports (:cpuid-field-target:FPTRANSA.cpuid.EXTENSIONS.FPTRANSA_ACCURACY.HEADER.MAX_INDEX:)\texttt{MAX\_INDEX}\texttt{=0x0044}.
A contract ID identifies the reference function, input domain, special-value behavior, error metric, exact anchors,
and mathematical properties; it is independent of instruction encoding placement. A different tuple of those
properties uses a different ID, and a retired ID is not reused. The contracts defined here use
(:ref:FPTRANSA.cpuid.EXTENSIONS.FPTRANSA_ACCURACY.ACCURACY.CONTRACT_REVISION:)\texttt{=1}. A changed guarantee, domain, or error metric names a different contract ID.
An implementation that advertises a smaller certified ULP bound strengthens the existing contract and retains the
same ID and revision.

The sparse ID space reserves \texttt{0x0001..0x000F} for reduced-domain circular trigonometric functions,
\texttt{0x0010..0x001F} for inverse trigonometric functions, \texttt{0x0020..0x002F} for hyperbolic and inverse
hyperbolic functions, \texttt{0x0030..0x003F} for exponential functions, and
\texttt{0x0040..0x004F} for logarithmic functions. The low-nibble-zero selector of each family is unassigned, and
actual contracts in each family begin at low nibble one.

\begin{BedrockListedFormatDiagram}{FPTRANSA Accuracy Result Format}
\BedrockFormatRowRange{result[63:0]}{63}{0}{%
\BedrockFormatField{P}{1}
\BedrockFormatReserved{(:term:base.terms.reserved:)}{23}
\BedrockFormatField{REV}{8}
\BedrockFormatField{D ULP}{16}
\BedrockFormatField{S ULP}{16}
}
\end{BedrockListedFormatDiagram}
{\small
\texttt{P} is (:ref:FPTRANSA.cpuid.EXTENSIONS.FPTRANSA_ACCURACY.ACCURACY.PRESENT:); \texttt{REV} is (:ref:FPTRANSA.cpuid.EXTENSIONS.FPTRANSA_ACCURACY.ACCURACY.CONTRACT_REVISION:); the D and S ULP
(:term:base.terms.field|plural:) use unsigned Q8.8.
\par}

\BedrockTableCaption{FPTRANSA Accuracy Result}
\begin{BedrockTabular}{@{}>{\raggedright\arraybackslash}p{0.78in}>{\raggedright\arraybackslash}p{1.20in}>{\raggedright\arraybackslash}p{3.55in}@{}}
\toprule
\rowcolor{BedrockHeaderFill}
\textbf{Bits} & \textbf{Field} & \textbf{Meaning}\\
\midrule
\texttt{0..15} & (:cpuid-field-target:FPTRANSA.cpuid.EXTENSIONS.FPTRANSA_ACCURACY.ACCURACY.S_MAX_ULP_Q8_8:)\texttt{S\_MAX\_ULP\_Q8\_8} & certified worst-case S-format ULP bound, rounded upward to unsigned Q8.8\\
\texttt{16..31} & (:cpuid-field-target:FPTRANSA.cpuid.EXTENSIONS.FPTRANSA_ACCURACY.ACCURACY.D_MAX_ULP_Q8_8:)\texttt{D\_MAX\_ULP\_Q8\_8} & certified worst-case D-format ULP bound, rounded upward to unsigned Q8.8\\
\texttt{32..39} & (:cpuid-field-target:FPTRANSA.cpuid.EXTENSIONS.FPTRANSA_ACCURACY.ACCURACY.CONTRACT_REVISION:)\texttt{CONTRACT\_REVISION} & value 1 for the contracts defined here\\
\texttt{40..62} & (:term:base.terms.reserved:) & zero\\
\texttt{63} & (:cpuid-field-target:FPTRANSA.cpuid.EXTENSIONS.FPTRANSA_ACCURACY.ACCURACY.PRESENT:)\texttt{PRESENT} & the instruction and both S and D encodings are available\\
\bottomrule
\end{BedrockTabular}

For a certified bound $K$, the (:term:base.terms.field:) value is $\lceil 256K\rceil$ and the advertised bound is the (:term:base.terms.field:) divided by
256. Thus 0.5, 1, 1.5, 2, and 4 ULP encode as \texttt{0x0080}, \texttt{0x0100}, \texttt{0x0180},
\texttt{0x0200}, and \texttt{0x0400}. A present contract has both (:term:base.terms.field|plural:) in \texttt{0x0001..0x0400}.
An unrecognized returned revision does not satisfy the contracts defined here.

Software first checks (:ref:FP.cpuid.EXTENSIONS.DIRECTORY.FEATURES.FP:) and (:ref:FPTRANSA.cpuid.EXTENSIONS.DIRECTORY.FEATURES.FPTRANSA:) in \texttt{EXTENSION\_DIRECTORY}, then queries the intended
contract ID in \texttt{FPTRANSA\_ACCURACY}. It uses the instruction only when (:ref:FPTRANSA.cpuid.EXTENSIONS.FPTRANSA_ACCURACY.ACCURACY.PRESENT:) is set, the returned
(:ref:FPTRANSA.cpuid.EXTENSIONS.FPTRANSA_ACCURACY.ACCURACY.CONTRACT_REVISION:) is 1, and the S or D bound required by the call site is no larger than its quality
threshold; otherwise it
does not use the instruction under that accuracy contract. For example, a D-format (:ref:FPTRANSA.instructions.FLOG2A:) path requiring at most 2 ULP accepts selector
\texttt{0x0000000100020043} only when the returned D (:term:base.terms.field:) is at most \texttt{0x0200}.

The selector is $(\texttt{0x00000001} \ll 32) \mathbin{|} (\texttt{0x0002} \ll 16) \mathbin{|}
\textit{contract-id}$. For example, (:ref:FPTRANSA.instructions.FSINA:) uses selector \texttt{0x0000000100020001}, and
(:ref:FPTRANSA.instructions.FLOG2A:) uses \texttt{0x0000000100020043}. An implementation that certifies (:ref:FPTRANSA.instructions.FSINA:) at 1.5 ULP for
S and 2 ULP for D returns \texttt{0x8000000102000180}.

(:ref:FPTRANSA.cpuid.EXTENSIONS.FPTRANSA_ACCURACY.ACCURACY.PRESENT:)\texttt{=1} establishes the accuracy contract for both S and D encodings but does not waive the
instruction\textquotesingle{}s (:term:base.terms.ordinary:) legality checks. The S and D bounds are worst-case bounds over every input in
the contract domain and apply to every execution on the (:term:base.terms.logical_processor:) that returned the CPUID result until reset.
Executing the instruction with (:ref:FPTRANSA.cpuid.EXTENSIONS.FPTRANSA_ACCURACY.ACCURACY.PRESENT:)\texttt{=0} raises (:ref:base.events.EXCEPTION.INVALID_OPERAND_RELATION:)
before operand or floating-point-state access.

\Needspace{1.25in}
\BedrockTableCaption{FPTRANSA Accuracy Contracts}
\begin{BedrockLongTable}{@{}>{\raggedright\arraybackslash}p{0.72in}>{\raggedright\arraybackslash}p{0.78in}>{\raggedright\arraybackslash}p{1.20in}>{\raggedright\arraybackslash}p{2.05in}>{\raggedright\arraybackslash}p{0.90in}@{}}
\toprule
\rowcolor{ManualHeaderFill}
\textbf{Contract ID} & \textbf{Mnemonic} & \textbf{Reference} & \textbf{Domain} & \textbf{ISA maximum}\\
\midrule
\endfirsthead
\multicolumn{5}{@{}l}{\scriptsize\itshape Table \theBedrockTable\ (continued)}\\
\toprule
\rowcolor{ManualHeaderFill}
\textbf{Contract ID} & \textbf{Mnemonic} & \textbf{Reference} & \textbf{Domain} & \textbf{ISA maximum}\\
\midrule
\endhead
\texttt{0x0001} & (:ref:FPTRANSA.instructions.FSINA:) & sin(x) in radians & finite x with abs(x) $\leq$ pi / 4 after DAZ & S $\leq$ 4 ULP; D $\leq$ 4 ULP\\
\texttt{0x0002} & (:ref:FPTRANSA.instructions.FCOSA:) & cos(x) in radians & finite x with abs(x) $\leq$ pi / 4 after DAZ & S $\leq$ 4 ULP; D $\leq$ 4 ULP\\
\texttt{0x0003} & (:ref:FPTRANSA.instructions.FTANA:) & tan(x) in radians & finite x with abs(x) $\leq$ pi / 4 after DAZ & S $\leq$ 4 ULP; D $\leq$ 4 ULP\\
\texttt{0x0004} & (:ref:FPTRANSA.instructions.FSINCOSA:) & paired sin(x) and cos(x) in radians & finite x with abs(x) $\leq$ pi / 4 after DAZ & S $\leq$ 4 ULP; D $\leq$ 4 ULP\\
\texttt{0x0011} & (:ref:FPTRANSA.instructions.FASINA:) & asin(x) in radians & finite x with abs(x) $\leq$ 1 after DAZ & S $\leq$ 4 ULP; D $\leq$ 4 ULP\\
\texttt{0x0012} & (:ref:FPTRANSA.instructions.FACOSA:) & acos(x) in radians & finite x with abs(x) $\leq$ 1 after DAZ & S $\leq$ 4 ULP; D $\leq$ 4 ULP\\
\texttt{0x0013} & (:ref:FPTRANSA.instructions.FATANA:) & atan(x) in radians & all finite values and signed infinity after DAZ & S $\leq$ 4 ULP; D $\leq$ 4 ULP\\
\texttt{0x0021} & (:ref:FPTRANSA.instructions.FSINHA:) & sinh(x) & all finite values and signed infinity after DAZ & S $\leq$ 4 ULP; D $\leq$ 4 ULP\\
\texttt{0x0022} & (:ref:FPTRANSA.instructions.FCOSHA:) & cosh(x) & all finite values and signed infinity after DAZ & S $\leq$ 4 ULP; D $\leq$ 4 ULP\\
\texttt{0x0023} & (:ref:FPTRANSA.instructions.FTANHA:) & tanh(x) & all finite values and signed infinity after DAZ & S $\leq$ 4 ULP; D $\leq$ 4 ULP\\
\texttt{0x0024} & (:ref:FPTRANSA.instructions.FATANHA:) & atanh(x) & finite x with abs(x) $\leq$ 1 after DAZ & S $\leq$ 4 ULP; D $\leq$ 4 ULP\\
\texttt{0x0031} & (:ref:FPTRANSA.instructions.FETOXA:) & $e^x$ & all finite values and signed infinity after DAZ & S $\leq$ 4 ULP; D $\leq$ 4 ULP\\
\texttt{0x0032} & (:ref:FPTRANSA.instructions.FETOXM1A:) & $e^x-1$ without an intermediate rounded exponential & all finite values and signed infinity after DAZ & S $\leq$ 4 ULP; D $\leq$ 4 ULP\\
\texttt{0x0033} & (:ref:FPTRANSA.instructions.FTWOTOXA:) & $2^x$ & all finite values and signed infinity after DAZ & S $\leq$ 4 ULP; D $\leq$ 4 ULP\\
\texttt{0x0034} & (:ref:FPTRANSA.instructions.FTENTOXA:) & $10^x$ & all finite values and signed infinity after DAZ & S $\leq$ 4 ULP; D $\leq$ 4 ULP\\
\texttt{0x0041} & (:ref:FPTRANSA.instructions.FLOGNA:) & ln(x) & positive finite values and positive infinity after DAZ; signed zero is the DZ boundary & S $\leq$ 4 ULP; D $\leq$ 4 ULP\\
\texttt{0x0042} & (:ref:FPTRANSA.instructions.FLOGNP1A:) & ln(1 + x) without an intermediate rounded addition & finite x $\geq -1$ and positive infinity after DAZ; $-1$ is the DZ boundary & S $\leq$ 4 ULP; D $\leq$ 4 ULP\\
\texttt{0x0043} & (:ref:FPTRANSA.instructions.FLOG2A:) & log2(x) & positive finite values and positive infinity after DAZ; signed zero is the DZ boundary & S $\leq$ 4 ULP; D $\leq$ 4 ULP\\
\texttt{0x0044} & (:ref:FPTRANSA.instructions.FLOG10A:) & log10(x) & positive finite values and positive infinity after DAZ; signed zero is the DZ boundary & S $\leq$ 4 ULP; D $\leq$ 4 ULP\\
\bottomrule
\end{BedrockLongTable}

